\clearpage
\fancyhead[L]{\small ELEKTRODYNAMIK.}
\thispagestyle{empty}
\pdfbookmark[1]{Elektrodynamik. Ausgabe und Quelle}{electrodynamics-edition}
\begin{center}\Large\textbf{Zur mathematischen Theorie der elektrodynamischen Wirkungen}\end{center}

Der Titel der Nachlassabteilung und die Abschnitte [1]--[21] von Carl Friedrich Gauss’ Arbeit in den \textit{Werken}, Band V (Göttingen, 1867), S.~601--630, liegen hier als KI-Übertragung in modernes Deutsch mit KI-erstelltem Satz vor. Die Abschnitte [1]--[3] behandeln die Stromintensität in einem Leitungsnetz, den Widerstand und die Intensitäten in einem bestimmten Stromkreis sowie die Wirkung zweier Stromelemente aufeinander. Die Abschnitte [4] und [5] behandeln die Umschlingungen zweier Linien und das Gesetz des galvanischen Stroms. Abschnitt [6] behandelt die Kraft zwischen Stromelementen und die elektromotorische Kraft. Die Abschnitte [7] und [8] enthalten das Induktionsgesetz und die Gesetze für zwei geschlossene Leiter. Abschnitt [9] behandelt die Induktion durch Bewegung. Die Abschnitte [10]--[12] geben Ausdrücke für die Wirkungen galvanischer Ströme an, behandeln geradlinige Polygone und setzen elektrische und magnetische Größen zueinander in Beziehung. Die Abschnitte [13] und [14] stellen das Grundgesetz für galvanische Ströme in zwei Fassungen dar. Die Abschnitte [15]--[17] behandeln die Kugelfläche, die elektromotorische Kraft bei der Entstehung eines Stroms und die Wirkung eines bewickelten Rings. Abschnitt [18] behandelt das Gleichgewicht und die Wirkung eines galvanischen Stroms auf einen Körper mit positiven und negativen Massen. Abschnitt [19] enthält Gauss’ Beweis von Ampères Fundamentalsatz. Abschnitt [20] ist Gauss’ Brief an Wilhelm Weber vom 19. März 1845. Abschnitt [21] stellt Beziehungen zwischen physikalischen Größen zusammen. Die Abbildungen sind in bearbeitbarem \LaTeX{} gezeichnet.

\medskip\noindent
\href{https://gdz.sub.uni-goettingen.de/id/PPN236006339}{Deutsche Werkausgabe (1867)}

Die Notiz über Umschlingungen von 1833 liegt auch in englischer Übersetzung bei Renzo L. Ricca und Bernardo Nipoti, \textit{Gauss’ Linking Number Revisited}, Abb.~1(b), S.~1327, sowie in der Übersetzung von Martin Mathieu bei Charles Nash, \textit{Topology and Physics---a Historical Essay}, §~1, S.~3--4, Anm.~3, vor. Wortlaut und Formel folgen hier der deutschen Werkausgabe.

\medskip\noindent
\href{https://webhomes.maths.ed.ac.uk/~v1ranick/papers/ricca.pdf#page=3}{Ricca und Nipoti (2011)}\quad
\href{https://arxiv.org/abs/hep-th/9709135v3}{Übersetzung von Mathieu bei Nash (1997)}
\par\medskip\noindent Clemens Schaefer gibt die Zeitangabe im datierten Vermerk zu Abschnitt [7] ausgeschrieben wieder, in \href{https://www.nature.com/articles/128339a0}{\textit{Gauss’s Investigations on Electrodynamics}} (1931), S.~341, Anm.~2.
\par\medskip\noindent Der \href{https://archive.org/details/Werkecarlf05gausrich/page/615/mode/1up}{zweite Abdruck der \textit{Werke}, Band~V (Göttingen, 1877), S.~615}, enthält die Zeichenübersicht zu Abschnitt [10] und die logarithmischen Ausdrücke in Abschnitt [11].
\par\medskip\noindent A.~K.~T. Assis hat Abschnitt [13] ins Englische übersetzt, in \href{https://www.ifi.unicamp.br/~assis/Weber-in-English-Vol-4.pdf#page=14}{\textit{Wilhelm Weber’s Main Works on Electrodynamics Translated into English}, Band~IV (Apeiron, 2021), Kap.~3, S.~11}. James Clerk Maxwell zitiert den ersten Satz auf Englisch in \href{https://en.wikisource.org/wiki/Page:A_Treatise_on_Electricity_and_Magnetism_-_Volume_2.djvu/461}{\textit{A Treatise on Electricity and Magnetism}, Band~II (1873), §~851, S.~429}.
\par\medskip\noindent Eine englische Übersetzung von Abschnitt [15] stammt von user2554 und steht in der \href{https://hsm.stackexchange.com/questions/7503/several-questions-about-gausss-contributions-to-electromagnetism}{Fassung vom 15. August 2025 einer Frage bei History of Science and Mathematics Stack Exchange}. Die Abbildungen zu den Abschnitten [16] und [17] stehen auch im \href{https://archive.org/details/Werkecarlf05gausrich/page/621/mode/1up}{zweiten Abdruck der \textit{Werke}, Band~V (1877), S.~621}. Maxwell führt das Ergebnis von Abschnitt [17] in \href{https://en.wikisource.org/wiki/Page:A_Treatise_on_Electricity_and_Magnetism_-_Volume_2.djvu/344}{\textit{A Treatise on Electricity and Magnetism}, Band~II (1873), §~706, S.~312} an.
\par\medskip\noindent Susan P. Johnsons englische Übersetzung des Briefes in Abschnitt [20] steht in \href{https://21sci-tech.com/articles/Atomic_Science.pdf#page=22}{\textit{Text of the Gauss-Weber 1845 Correspondence}}, herausgegeben von Laurence Hecht, \textit{21st Century Science \& Technology} 9(3), Herbst 1996, S.~42--43. Hecht verfasste den begleitenden Aufsatz. Die Anmerkung der Redaktion auf S.~41 nennt Karl Krause und Alexander Hartmann als Transkribenten der drei benachbarten Briefe von Weber an Gauss.
\par\smallskip\noindent In Gauss’ Brief bedeutet \textit{in verbis facilior} „nachgiebiger in Fragen des Wortlauts“ und \textit{Nil actum reputans si quid superesset agendum} „nichts als getan ansehen, solange noch etwas zu tun blieb“.
\par\smallskip\noindent Abschnitt [21] steht auch im \href{https://archive.org/download/Werkecarlf05gausrich/page/n641}{zweiten Abdruck der \textit{Werke}, Band~V (1877), S.~630}.
\clearpage
